\documentclass{article}
\usepackage[T1]{fontenc}
\usepackage{lmodern}
\usepackage{graphicx}
\usepackage{amsmath,amssymb}
\usepackage[a4paper,margin=2cm]{geometry}
\usepackage[absolute,overlay]{textpos}
\setlength{\TPHorizModule}{1mm}
\setlength{\TPVertModule}{1mm}
\usepackage[indonesian]{babel}
\usepackage{hyperref}
\setlength{\emergencystretch}{2em}
% unit: Halaman 11

\begin{document}

% === Halaman 11 (llm) ===

\begin{itemize}
    \item Contoh di dalam Caesar Cipher, misalkan penyerang tahu cipherteks:
    
    KHOOR ZRUOG
    
    dan juga tahu plainteks-nya seharusnya mengandung kata:
    
    HELLO
    
    Dari pasangan $\text{H} \rightarrow \text{K}$, penyerang tahu pergeseran $(k) = 3$.
\end{itemize}

Dengan kunci $k = 3$, semua cipherteks lain bisa didekripsi.

\end{document}
